
\begin{frame}
\frametitle{Le partitionnement (\textit{sharding})}

Le partitionnement consiste à  

\begin{itemize}
  \item \alert{découper} une (grande) collection
      en \alert{fragments} en fonction d'une \alert{clé}\,;
      \item placer chaque fragment sur un \alert{serveur}\,;
      \item Maintenir un \alert{répertoire} indiquant  que telle \alert{clé} se trouve dans tel \alert{fragment}
          sur tel \alert{serveur}
\end{itemize}
 
\vfill

Le partitionnement apporte la \alert{répartition de charge} (\textit{load balancing})

\vfill

Il doit être \alert{dynamique} (ajout/retrait de serveurs) pour s'adapter ``élastiquement'' 
 
 \vfill
 
 S'applique à des collections de paires (clé, valeur), où valeur est n'importe quelle information structurée.
\end{frame}

\begin{frame}
\frametitle{Opérations}

Les opérations de type ``dictionnaire'' peuvent être transmises à \alert{un seul} serveur\,:

\begin{itemize}
   \item $get(k): v$, recherche par clé
   \item $put (k, v)$, insertion
   \item $delete(k)$, suppression
   \item $range(k_1, k_2)$, recherche par intervalle (peut impliquer plusieurs serveurs, ne marche pas avec le hachage)
\end{itemize}

\vfill

\alert{Pour toutes les autres opérations}\,: tous les serveurs effectuent l'opération localement.

\vfill

Exemple\,: recherche sur une clé secondaire, s'effectue sur \alert{tous} les serveurs, avec si possible des index locaux.
\end{frame}

\section{Par tri}

\begin{frame}
\frametitle{Première étape\,: il faut trier}

Il faut disposer d'une structure triant les documents sur la clé. Deux possibilités\,:

\vfill

\alert{1) On trie la collection (clé, document)}

\begin{itemize}
  \item indexation dite ``plaçante'': l'emplacement d'un document est déterminé par sa clé.
  \item efficace en recherche (pas d'indirection)
  \item peu efficace en mise à jour (nécessite des déplacements)  
  \item on peut créer un seul index plaçant par collection
\end{itemize}

\alert{2) On trie une structure indépendante constituée des paires (clé, adresse) }

\begin{itemize}
  \item indexation non-plaçante\,: les documents peuvent être insérés n'importe où (efficace)
  \item un peu moins efficace en recherche (on trouve une adresse dans l'index, pas le document lui-même)
  \item autant d'index que l'on veut
\end{itemize}

\end{frame}



\begin{frame}
\frametitle{Seconde étape (optionelle)\,: on crée des niveaux d'index}

Illustration avec \alert{collection triée}, la clé est le titre.

\figSlideWithSize{idxnondense}{10cm}

On a découpé en \alert{4 fragments}, 
indexés avec un niveau.

\end{frame}


\begin{frame}
\frametitle{Quelques calculs, cas d'une collection triée}

On a un million de vidéos, soit $2^{20}$, \alert{triés} sur la clé.

\vfill

\alert{Sans index}
\begin{itemize}
  \item Recherche par dichotomie, soit 20 accès disques pour trouver une vidéo
  \item Temps\,: 20 fois 10 ms = 200 ms (au pire du pire)
\end{itemize} 

\alert{Avec index}
\begin{itemize}
  \item Mon index comprend 1 million de paires (clé, adresse)
  \item Supposons qu'une adresse occupe 8 octets, la clé 8 octets aussi\,: l'index occupe 16 MO 
  \item Temps\,: recherche quasi instantanée dans l'index, 1 accès au disque (10 ms au pire) 
\end{itemize} 

\begin{noteblock}{Conclusion}
L'index est tout petit (tient en mémoire), permet des accès très rapides même sur de très grandes collections.
\end{noteblock}
\end{frame}


\begin{frame}
\frametitle{A quoi sert de \alert{trier} la collection\,?}

Différence (subtile)\,: si la collection est triée, je peux créer
un \alert{index non-dense}, sinon l'index \alert{doit} être dense.

\vfill
\alert{Index non-dense}

\begin{itemize}
  \item Je découpe ma collection en \alert{fragments} (quelques MO ou dizaines de MO)
  \item Chaque fragment couvre un intervalle
  \item \alert{J'indexe seulement les intervalles}
\end{itemize} 

Exemple dans le cas précédent: je découpe ma collection en fragments de 4 films; j'indexe
250\,000 fragments (1M/4).

\vfill

\alert{Si la collection n'est pas triée, je dois créer un index dense} (pas le choix)
\begin{itemize}
  \item Chaque document est représenté dans l'index
\end{itemize} 

\end{frame}

\begin{frame}
\frametitle{Exemple d'un index dense}

On indexe sur l'année, en y mettant \alert{toutes} les clés.

\figSlideWithSize{idxdense2}{10cm}

Presque aussi efficace, sauf pour les recherches par intervalle.
\end{frame}


\section{Par hachage}


\begin{frame}
\frametitle{A la base\,: table de hachage et fichiers}

On \textit{calcule} la position d'un enregistrement
d'après la clé.
\begin{itemize}
  \item un \textit{fonction de hachage} $h$ associe des
         valeurs de clé à des adresses de bloc
   \item $h$ doit répartir uniformément les enregistrements
            dans les $n$ blocs alloués à la structure
  \item recherche d'un enregistrement =$>$ calcul de l'endroit
          où il se trouve.
\end{itemize}
Simple et efficace\,!
\end{frame}


%------------------------------------------------------------
\begin{frame}{Exemple}

On veut créer une structure de hachage pour nos 16 films
(hyp.\,: 4 enregistrements   par page)
\begin{itemize}
  \item on alloue 5 pages (pour garder une marge de man{\oe}uvre)
  \item un répertoire à 5 entrées (0 à 4) pointe vers les pages
   \item On définit la fonction $h(titre) = rang(titre[0])\, mod\, 5$
\end{itemize}

\end{frame}

%------------------------------------------------------------
\begin{frame}{Le résultat}

  \figSlideWithSize{hachage}{10cm}

\end{frame}


\begin{frame}{Recherches}

\begin{itemize}
  \item Par clé\,? Oui:\\
     \texttt{SELECT * FROM  Film WHERE titre = 'Impitoyable'}

   \item Par préfixe\,? Ici, oui.\\
      \texttt{SELECT * FROM  Film WHERE titre LIKE 'M\%'}

   \item Par intervalle\,: non\,!\\

\texttt{SELECT *
FROM  Film
WHERE titre BETWEEN 'Annie Hall' AND 'Easy Rider'}
  
\end{itemize}

\end{frame}
%------------------------------------------------------------
\begin{frame}{Mises à jour\,: ça se gâte}

La structure simple décrite précédemment n'est pas
\textbf{dynamique}
\begin{itemize}
  \item On ne peut pas changer un enregistrement de place

   \item Donc  il faut créer un chaînage de pages quand
             une page déborde

  \item Et donc les performances se dégradent...
\end{itemize}

\end{frame}
%------------------------------------------------------------
\begin{frame}{Exemple\,: insertion de Citizen Kane}

  \figSlideWithSize{hachage-2}{10cm}

\end{frame}

%------------------------------------------------------------
\begin{frame}{Hachage dynamique}

Objectif\,: réorganiser la table de hachage en fonction
des insertions et suppressions. 

\begin{itemize}
  \item le nombre d'entrées dans le répertoire est une puissance
          de 2
  \item la fonction $h$ donne toujours un entier sur 4 octets (32 bits)  
\end{itemize}
Idée de base\,: 
 on utilise les $n$ premiers bits du résultat de la fonction, avec
$n < 32$

\end{frame}
%------------------------------------------------------------
\begin{frame}{Exemple\,: le hachage des 16 films}

\begin{center}
\begin{tabular}{l|c}
titre & $h$(titre)   \\
\hline
Vertigo     & 01110010  \\
Brazil      & 10100101  \\
Twin Peaks  & 11001011  \\
Underground & 01001001  \\
Easy Rider  & 00100110 \\
Psychose    &  01110011 \\
Greystoke   &  10111001 \\
Shining     &  11010011 \\
\hline
\end{tabular}
\end{center}

\end{frame}
%------------------------------------------------------------
\begin{frame}{Construction de la table}

On départ on utilise seulement le premier bit
de la fonction
\begin{itemize}
  \item Deux valeurs possibles\,: 0 et 1
  \item Donc deux entrées, et deux blocs
  \item L'affectation d'un enregistrement dépend
        du premier bit de sa fonction de hachage  
\end{itemize}

=$>$ pour l'instant on reste dans un cadre classique

\end{frame}
%------------------------------------------------------------
\begin{frame}{Avec 5 films}

  \figSlide{hashext-1}

\end{frame}
%------------------------------------------------------------
\begin{frame}{Insertions}

Supposons 3 films par bloc. L'insertion de Psychose
(valeur 01110011) entraîne le débordement du premier bloc.

\begin{itemize}
  \item On double la taille du répertoire
  \item On alloue un nouveau bloc pour l'entrée 01
  \item Les entrées 10 et 11 pointent sur le \textit{même} bloc.
\end{itemize}
On agrandit seulement ce qui est nécessaire.

\end{frame}
%------------------------------------------------------------
\begin{frame}{Illustration}

  \figSlide{hashext-2}

\end{frame}
%------------------------------------------------------------
\begin{frame}{Insertions suivantes}

Plusieurs cas
\begin{itemize}
  \item on insère dans un bloc plein, mais plusieurs
           entrées pointent dessus
  \begin{itemize}
    \item on alloue un nouveau bloc, et on répartit les pointeurs
  \end{itemize}
  \item on insère dans un bloc plein, associé à une entrée
  \begin{itemize}
    \item on double à nouveau le nombre d'entrées
  \end{itemize}
Problème\,: le répertoire peut devenir très grand.
\end{itemize}

\end{frame}
%------------------------------------------------------------
\begin{frame}{Greystoke 10111001 et Shining 11010011}

  \figSlide{hashext-3}

\end{frame}




\section{Et en distribué?}


\begin{frame}
\frametitle{Architecture d'un système avec partitionnement}

\figSlide{archi-partition}
      
\begin{itemize}
  \item Routeur = \textit{Master}, \textit{Balancer}, \textit{Config server}, ...
  \item Fragment = \textit{chunk}, \textit{tablet}, \textit{region}, \textit{bucket}, ...
\end{itemize}

\begin{remark}
Chaque serveur gère une réplication locale pour tolérance aux pannes.
\end{remark}
\end{frame}


\begin{frame}
\frametitle{Techniques de partitionnement}

\begin{itemize}
   \item par \alert{intervalle}\,: documents triés sur la clé, puis groupés par intervalles.
   
     Exemples représentatifs\,: \mongodb/, HBase (BigTable).
     
   \item par \alert{hachage} sur la clé, avec répertoire de hachage.
   
       Exemples\,: \textit{memcached}, \textit{chord}, \textit{Dynamo/S3/Voldemort}, beaucoup d'autres.
 \end{itemize}
 

\begin{remark}
Se référer aux techniques d'indexation par arbre-B et par hachage dans des environnements centralisés.
\end{remark}
\end{frame}



\begin{frame}
\frametitle{Notre application}

On veut gérer (légalement) une collection de vidéos.

\begin{itemize}
  \item chaque vidéo occupe en moyenne 500 MO
  \item on associe chaque vidéo à quelques méta données (titre, auteur, année, etc.)
  \item 1\,000 vidéos = 500 GO (1/2 TO), pas de souci avec un (bon) serveur.
  \item 1\,000\,000 vidéos = 500 TO (1/2 PO), il faut distribuer.
\end{itemize}

\begin{remark}
Dès qu'on parle de système distribué, il faut envisager le coût en stockage de la réplication. Multiplier
par 2 ou 3...
\end{remark}


\begin{noteblock}{Besoins}
Chercher une vidéo par son titre, les vidéos par année, intervalle d'années, etc.
\end{noteblock}

\end{frame}


\begin{frame}
\frametitle{Pour illustrer}

Un tout petit échantillon.

\begin{center}
\begin{small}
\begin{tabular}{r|l|c|c||r|l|c|c}
id &titre & année  & mp4  & id & titre & année  & mp4  \\
\hline
1 & Vertigo  & 1958 & ... & 9 & Annie Hall & 1977 & ... \\
2 & Brazil & 1984 & ... & 10&Jurasic Park & 1992 & ... \\
3&Twin Peaks & 1990 & ...&  11&Metropolis & 1926 & ... \\
4&Underground & 1995 & ... & 12&Manhattan & 1979 & ... \\
5&Easy Rider & 1969 & ... & 13&Reservoir Dogs & 1992 & ... \\
6&Psychose & 1960 & ... & 14&Impitoyable & 1992 & ... \\
7&Greystoke & 1984 & ... & 15&Casablanca & 1942 & ... \\
8&Shining & 1980 & ... & 16&Smoke & 1995 & ... \\
\hline
\end{tabular}
\end{small}
\end{center}

Quelle clé considère-t-on\,? On va y revenir.

\end{frame}